\begin{table}[tb]
\centering
\caption{Label-level comparisons on the three main backbones: mean difference in held-out AUROC (classification) or Dice (segmentation), its 95\% interval over cohorts, labels and held-out cases, and the number of cohorts (datasets) with a positive mean difference. Each label counts once with its per-task and joint differences averaged, except in the two LoRA $-$ Dec sub-rows. LoRA $-$ Attn compares two methods with the same attentive head. Pooled estimates (``all'') use TAP-CT, SPECTRE and Pillar-0; VoCo is in Table~\ref{tab:keyvoco}. $^{\dagger}$ = did not train with the default setting (Section~\ref{sec:fullft}).}
\label{tab:key}
\adjustbox{max width=\textwidth}{\small\begin{tabular}{@{}lccccc@{}}
\toprule
Comparison & TAP-CT & SPECTRE & Pillar-0 & All & Cohorts $>0$ \\
\midrule
\multicolumn{6}{@{}l}{\textit{Classification (AUROC)}} \\
LoRA $-$ Lin & \mci{+0.032}{+0.007}{+0.053} & \mci{+0.037}{+0.009}{+0.075} & \mci{+0.033}{+0.021}{+0.044} & \mci{+0.034}{+0.017}{+0.052} & 9/10 \\
LoRA $-$ Attn & \mci{+0.022}{+0.006}{+0.042} & \mci{+0.033}{+0.003}{+0.078} & \mci{+0.034}{+0.021}{+0.052} & \mci{+0.027}{+0.012}{+0.048} & 10/10 \\
Attn-only FT $-$ LoRA & \mci{+0.019}{-0.001}{+0.054} & \mci{+0.002}{-0.037}{+0.026} & \mci{+0.001}{-0.002}{+0.004} & \mci{+0.011}{-0.007}{+0.029} & 7/10 \\
Full FT $-$ Lin & \mci{+0.022}{-0.061}{+0.068} & \mci{-0.117}{-0.195}{-0.059}$^{\dagger}$ & \mci{+0.004}{-0.013}{+0.021} & \mci{-0.038}{-0.103}{+0.019} & 5/10 \\
joint $-$ per-task & \mci{+0.012}{-0.014}{+0.038} & \mci{-0.003}{-0.044}{+0.051} & \mci{-0.008}{-0.013}{-0.001} & \mci{+0.004}{-0.021}{+0.030} & 6/10 \\
\addlinespace
\multicolumn{6}{@{}l}{\textit{Segmentation (Dice)}} \\
Dec $-$ Lin & \mci{+0.085}{+0.055}{+0.109} & \mci{+0.061}{+0.027}{+0.088} & -- & \mci{+0.073}{+0.049}{+0.092} & 10/10 \\
LoRA $-$ Dec & \mci{-0.005}{-0.021}{+0.024} & \mci{+0.023}{-0.018}{+0.076} & -- & \mci{+0.009}{-0.015}{+0.046} & 8/10 \\
LoRA $-$ Dec, per-task & \mci{-0.021}{-0.036}{+0.002} & \mci{+0.006}{-0.003}{+0.031} & -- & \mci{-0.007}{-0.022}{+0.009} & 7/10 \\
LoRA $-$ Dec, joint & \mci{+0.011}{-0.021}{+0.066} & \mci{+0.040}{-0.039}{+0.147} & -- & \mci{+0.026}{-0.021}{+0.093} & 8/10 \\
Full FT $-$ Dec & \mci{0.000}{-0.024}{+0.016} & \mci{-0.003}{-0.035}{+0.009} & -- & \mci{-0.002}{-0.018}{+0.009} & 6/10 \\
joint $-$ per-task & \mci{-0.046}{-0.090}{-0.008} & \mci{-0.058}{-0.093}{-0.009} & -- & \mci{-0.052}{-0.082}{-0.022} & 0/10 \\
\bottomrule
\end{tabular}
}
\end{table}

We report held-out results of \numRunsMainN{} main runs and \numRunsRevisionN{} repeats and controls. Table~\ref{tab:key} summarises the main comparisons (per backbone in Figure~\ref{fig:key}, Appendix~\ref{app:perlabel}); per-cohort tables, ranks, per-label results and every paired test are in Appendix~\ref{app:perlabel}. The selection-split metric recomputed from the stored per-case predictions matched the value logged during training for every run (classification to $10^{-7}$; segmentation within $2\times10^{-3}$ for all but one of the \numSelcheckSegRunsN{} runs, whose per-case Dice was re-inferred after training).

\subsection{Classification}
\label{sec:cls}

\paragraph{Adapting the encoder improves on the frozen probes.} LoRA raises held-out AUROC over the linear probe by $\numCiClsloralinAllMeandiff$ at label level (95\% interval $\numCiClsloralinAllLodiff$ to $\numCiClsloralinAllHidiff$; positive in \numCiClsloralinAllPosN{} of \numCiClsloralinAllClustersN{} cohort--backbone pairs), with intervals above zero on each backbone: $\numCiClsloralinTapctMeandiff$ on TAP-CT, $\numCiClsloralinSpectreMeandiff$ on SPECTRE and $\numCiClsloralinPillarZMeandiff$ on Pillar-0. LoRA also changes the head, from linear to attentive, but the frozen attentive probe alone adds only $\numCiClsattnlinAllMeandiff$ ($\numCiClsattnlinAllLodiff$ to $\numCiClsattnlinAllHidiff$), and against this matched-head baseline LoRA gains $\numCiClsloraattnAllMeandiff$ ($\numCiClsloraattnAllLodiff$ to $\numCiClsloraattnAllHidiff$; positive in \numCiClsloraattnAllPosN{} of \numCiClsloraattnAllClustersN{} cohort--backbone pairs), so the gain comes from adapting the encoder. Leaving out FDG-PET/CT and COVID-CTset, two sources in which we did not rule out shortcuts, leaves it unchanged ($\numCiClsloralinnoscAllMeandiff$, $\numCiClsloralinnoscAllLodiff$ to $\numCiClsloralinnoscAllHidiff$). Per label--mode pair, LoRA is never significantly behind the linear probe and significantly ahead in \numPairedTapctLinLoraSigbetter\%, \numPairedSpectreLinLoraSigbetter\% and \numPairedPillarZLinLoraSigbetter\% of the pairs on TAP-CT, SPECTRE and Pillar-0; significance follows the held-out size rather than the effect (mean gain $\numSigbysizeSmallMeandiff$ and \numSigbysizeSmallSig\% significant below 100 held-out cases, $\numSigbysizeLargeMeandiff$ and \numSigbysizeLargeSig\% above 400; Figure~\ref{fig:perlabel}). LoRA is ahead of the linear probe in \numSeedsclsLoraLinAllPosN{} of \numSeedsclsLoraLinAllPairsN{} cohort--seed pairs of the seed repeats, and the gain is largest where the probe is weakest ($\numGainbinLowMeandiff$ for labels whose probe reaches a selection AUROC of at most 0.6, $\numGainbinTopMeandiff$ above 0.8).

\paragraph{Training the attention weights in full has the best mean rank, but its gain over LoRA is uncertain.} At the default settings and over the \numRankSlotsN{} label--mode pairs, attention-only fine-tuning has the best mean rank (\numRankAttnftMean{}; LoRA \numRankLoraMean{}, LoRA+ABMIL \numRankBabmilMean{}, DoRA \numRankDoraMean{}, linear probe \numRankLinMean{}; Figure~\ref{fig:ranks}). Its label-level difference to LoRA is $\numCiClsattnftloraAllMeandiff$ ($\numCiClsattnftloraAllLodiff$ to $\numCiClsattnftloraAllHidiff$): $\numCiClsattnftloraTapctMeandiff$ on TAP-CT but $\numCiClsattnftloraSpectreMeandiff$ on SPECTRE and $\numCiClsattnftloraPillarZMeandiff$ on Pillar-0. The LoRA variants do not improve on LoRA: on the same pairs, LoRA$_{\text{late}}$ changes AUROC by $\numVsloraAllLoragMeandiff$, DoRA by $\numVsloraAllDoraMeandiff$, rsLoRA by $\numVsloraAllRsloraMeandiff$, LoRA+ABMIL by $\numVsloraAllBabmilMeandiff$, and the two-stage schedules by $\numVsloraAllLinThenLoragMeandiff$ (Lin$\to$LoRA) and $\numVsloraAllAttnThenLoragMeandiff$ (Attn$\to$LoRA). The rsLoRA deficit on TAP-CT is a learning-rate effect: at a quarter of the rate, which matches its update scale to LoRA, it gains $\numCtrlTapctRsloraLowlrRsloraMeandiff$ (Appendix~\ref{app:revision}).

\paragraph{Full fine-tuning did not train on SPECTRE with the default setting.}\label{sec:fullft} The held-out probabilities of SPECTRE Full FT barely vary between cases (median standard deviation \numSpreadSpectreFullft{}, against at least \numSpreadOthermin{} for every other method on the three backbones), and it loses $\numCiClsfullftlinSpectreMeandiff$ to the linear probe. On TAP-CT it changes AUROC by $\numCiClsfullftlinTapctMeandiff$ ($\numCiClsfullftlinTapctLodiff$ to $\numCiClsfullftlinTapctHidiff$), with one collapsed joint run (FDG-PET/CT, \numClsAurocFdgTapctJointFullft{}), and it is the least stable method across seeds (median between-seed change of the cohort mean \numSeedsclsFullftSdmedian{}, LoRA \numSeedsclsLoraSdmedian{}). Of three alternative settings, layer-wise learning-rate decay (0.75 per block) gave SPECTRE Full FT the highest selection AUROC (\numFixSpectreFullftLlrdVal{}, against \numFixSpectreFullftLowlrVal{} with a lower learning rate and \numFixSpectreFullftAccumVal{} with gradient accumulation and clipping, which trained on \numFixSpectreFullftAccumTrainedN{} of \numFixSpectreFullftAccumCohortsN{} cohorts); with it Full FT is ahead of LoRA by $\numFixselSpectreLoraMeandiff$ on SPECTRE and level with it on TAP-CT ($\numFixselTapctLoraMeandiff$; Table~\ref{tab:revisioncls}), a comparison that favours Full FT, since LoRA was not re-tuned.

\paragraph{One joint model per cohort is as accurate as one model per label.} Joint minus per-task AUROC is $\numCiClsjointAllMeandiff$ ($\numCiClsjointAllLodiff$ to $\numCiClsjointAllHidiff$), at $1/K$ of the training steps of the $K$ per-task models; the sign varies by cohort from $\numJointdeltaMin$ to $\numJointdeltaMax$ (Figure~\ref{fig:joint}), and on Pillar-0 the joint model is behind on average, with an interval below zero ($\numCiClsjointPillarZMeandiff$, $\numCiClsjointPillarZLodiff$ to $\numCiClsjointPillarZHidiff$), though not for every method (Table~\ref{tab:cls}).

\subsection{Segmentation}
\label{sec:seg}

\paragraph{A decoder on the frozen encoder gives most of the gain.} Replacing the linear probe by the frozen decoder raises Dice by $\numCiSegdeclinAllMeandiff$ ($\numCiSegdeclinAllLodiff$ to $\numCiSegdeclinAllHidiff$), on every dataset of both backbones and significantly on \numSegpairedTapctLinDecSigbetter\% (TAP-CT) and \numSegpairedSpectreLinDecSigbetter\% (SPECTRE) of the structure--mode pairs. That is all of LoRA's gain over the linear probe on TAP-CT ($\numSegdeltaTapctLinLoraMeandiff$) and most of it on SPECTRE ($\numSegdeltaSpectreLinLoraMeandiff$).

\paragraph{Adapting the encoder on top of the decoder changes Dice little on average.} LoRA changes Dice relative to the frozen decoder by $\numCiSegloradecAllMeandiff$ ($\numCiSegloradecAllLodiff$ to $\numCiSegloradecAllHidiff$; $\numCiSegloradecTapctMeandiff$ on TAP-CT, $\numCiSegloradecSpectreMeandiff$ on SPECTRE). The sign depends on mode and dataset (Figure~\ref{fig:perstructure}): the largest gain is the SPECTRE AMOS22 joint model, where the decoder (\numSegAmosBBSpectreJointDec{}) falls below the linear probe (\numSegAmosBBSpectreJointLin{}) and LoRA reaches \numSegAmosBBSpectreJointLora{}, while on the TotalSegmentator joint model LoRA loses on both backbones (\numSegTotalsegSpectreJointLora{} against \numSegTotalsegSpectreJointDec{} on SPECTRE). LoRA is significantly ahead of the decoder on \numSegpairedSpectreDecLoraSigbetter\% of SPECTRE pairs and behind on \numSegpairedSpectreDecLoraSigworse\%; on TAP-CT on \numSegpairedTapctDecLoraSigbetter\% and \numSegpairedTapctDecLoraSigworse\%. Full fine-tuning has the best mean rank (\numSegrankFullftMean{}; LoRA \numSegrankLoraMean{}, Dec \numSegrankDecMean{}), on the strength of TAP-CT, but its mean change relative to the decoder is $\numCiSegfullftdecAllMeandiff$ ($\numCiSegfullftdecAllLodiff$ to $\numCiSegfullftdecAllHidiff$) because \numSegdeltaTapctDecFullftDropN{} TAP-CT structure--mode pairs lose more than 0.1 Dice. DoRA, LoRA$_{\text{late}}$ and Dec$\to$LoRA change Dice relative to the decoder by $\numCiSegdoradecAllMeandiff$, $\numCiSegloragdecAllMeandiff$ and $\numCiSegdecthenloradecAllMeandiff$, each with an interval that includes zero.

\paragraph{One model per structure is more accurate than one joint model at our budgets.} One joint model per dataset changes Dice by $\numCiSegjointAllMeandiff$ ($\numCiSegjointAllLodiff$ to $\numCiSegjointAllHidiff$), on none of the \numCiSegjointAllClustersN{} dataset--backbone pairs positive, most on the 15-structure AMOS22 ($\numSegjointdsAmosBBTapctMeandiff$, $\numSegjointdsAmosBBSpectreMeandiff$) and least on LiTS ($\numSegjointdsLitsTapctMeandiff$, $\numSegjointdsLitsSpectreMeandiff$). Much of the cost comes from small structures the joint model loses entirely: in \numSegcollapseN{} of \numSegcollapseCellsN{} structure--method pairs, all on the adrenal glands and the esophagus of AMOS22 and the merged glands of TotalSegmentator, the joint model scores below 0.01 Dice where the per-structure model reaches \numSegcollapsePertaskminValue{} to \numSegcollapsePertaskmaxValue{}; without them the mean difference over structure--method pairs moves from $\numSegjointTapctMeandiff$ to $\numSegjointTapctNocollapseMeandiff$ on TAP-CT and from $\numSegjointSpectreMeandiff$ to $\numSegjointSpectreNocollapseMeandiff$ on SPECTRE. Tripling the steps of the 15-structure joint models gained $\numSegbudgetgainTapctMeandiff$ and $\numSegbudgetgainSpectreMeandiff$, and at matched total compute (per-structure models of 1600 steps) the joint model is ahead with LoRA ($\numCmAmosBBTapctLoraJointdiff$ to $\numCmAmosBBSpectreLoraJointdiff$ on AMOS22) but behind with the frozen decoder and full fine-tuning in three of the four dataset--backbone pairs (Table~\ref{tab:computematched}). The gain from tripling the steps shows that the segmentation budget was short; \numConvSegAllLateBetter\% of segmentation runs also selected one of their last two validations, against \numConvClsAllLateBetter\% in classification, a share that the cosine decay of the learning rate inflates and that therefore indicates rather than measures under-training. The segmentation comparisons rank methods under a fixed budget, not at convergence.

\subsection{A fourth backbone: the hierarchical VoCo encoder}
\label{sec:voco}

VoCo differs from the three main backbones in its token interface (no class token, one stage-pooled token per $96^3$ window) and ran three of the five classification cohorts (\numLabelsVocoN{} labels) with the settings of the other backbones, so we report it separately (Appendix~\ref{app:voco}). In segmentation the findings hold: the decoder adds $\numCiSegdeclinVocoMeandiff$ Dice to a linear probe on the stride-32 output that segments almost nothing, LoRA changes Dice relative to the decoder by $\numCiSegloradecVocoMeandiff$ ($\numCiSegloradecVocoLodiff$ to $\numCiSegloradecVocoHidiff$), one joint model costs $\numCiSegjointVocoMeandiff$, and full fine-tuning is ahead of the decoder by $\numCiSegfullftdecVocoMeandiff$ ($\numCiSegfullftdecVocoLodiff$ to $\numCiSegfullftdecVocoHidiff$). In classification the picture differs: LoRA changes AUROC relative to the linear probe by $\numCiClsloralinVocoMeandiff$ ($\numCiClsloralinVocoLodiff$ to $\numCiClsloralinVocoHidiff$), behind on all \numCiClsloralinVocoClustersN{} cohorts, but more than half of this deficit comes from the head: the frozen attentive probe is itself behind the linear probe ($\numCiClsattnlinVocoMeandiff$), and against this matched-head baseline every method that adapts the encoder is level, with intervals including zero (LoRA $\numCiClsloraattnVocoMeandiff$, attention-only fine-tuning $\numCiClsattnftattnVocoMeandiff$, full fine-tuning $\numCiClsfullftattnVocoMeandiff$; Figure~\ref{fig:voco}). The only adapted method level with the linear probe is Lin$\to$LoRA, which keeps the linear head ($\numCiClslinthenloraglinVocoMeandiff$, $\numCiClslinthenloraglinVocoLodiff$ to $\numCiClslinthenloraglinVocoHidiff$). On VoCo, adapting the encoder therefore did not improve on a frozen probe with the same head, unlike on the three main backbones, and the attentive head on VoCo's 2688-dimensional window tokens was worse than a linear one. One joint model is again level with per-task models ($\numCiClsjointVocoMeandiff$). VoCo was run with the settings of the other backbones; whether a learning rate, head or window token chosen for it would change this we did not test, so the result bounds the transfer of the classification finding at the default settings rather than explaining it.
